% GENERATED FILE. Edit canonical lessons, then run npm run generate:books.
% canonical-source-hash: fnv1a64:45d0024a387d21c4
% canonical-lessons: BN-C416-kisti, BN-C416-dena, BN-C416-paona, BN-C416-shodh-kora, BN-C416-metano

\chapter{Debts and Instalments}
\label{ch:bn-debts-and-instalments}
\hlchaptermodality{416}

\begin{chapteropening}
\textbf{By the end of this chapter:} \emph{I can pay in instalments, say what I owe and am owed, and settle a bill.}

Complete the last lesson of chapter 416: I can pay in instalments, say what I owe and am owed, and settle a bill.
\end{chapteropening}

\section[\texorpdfstring{\bn{কিস্তি} (kisti)}{কিস্তি (kisti)}]{\bn{কিস্তি} (kisti) — an instalment}

\label{lesson:BN-C416-kisti}



\begin{quote}
Before the new one: say the Bengali for to invest, then the Bengali for to lose (money or property).
\end{quote}

\subsection*{You'll want to know: \bn{কিস্তি}}
\textbf{\bn{কিস্তি}} — \emph{kisti} — \textquotedblleft{}an instalment\textquotedblright{}.

\begin{quote}
\bn{অরুণ} \bn{সাইকেলটা} \bn{কিস্তিতে} \bn{কিনল।}
\end{quote}

\begin{quote}
\emph{ôrun sāikelṭā kistite kinlô.}
\end{quote}

\begin{quote}
Arun bought the bicycle in instalments.
\end{quote}

\begin{cousinweb}
\textbf{\bn{কিস্তি}} (\emph{kisti}) is usually traced to Arabic \emph{qisṭ}, \textquotedblleft{}a share, a portion\textquotedblright{}, through Persian \emph{qist}, \textquotedblleft{}an instalment\textquotedblright{}.
\end{cousinweb}

\begin{grammarlens}[title={What you've built}]
One more word for getting around.
\end{grammarlens}

\subsection*{Your turn}
\begin{itemize}
\raggedright
  \item \emph{Say it:} \emph{kisti}
  \item \emph{Say it:} \emph{kisti}, once more
  \item \emph{Say it:} \emph{ôrun sāikelṭā kistite kinlô.} — Arun bought the bicycle in instalments.
  \item \emph{Recall:} say the Bengali for to invest, then the Bengali for to lose (money or property), then say \emph{kisti} again
\end{itemize}

\begin{tcolorbox}[breakable,colback=teal!4,colframe=teal!35!black,title={Before you move on}]
What does \bn{কিস্তি} mean? (\textquotedblleft{}An instalment\textquotedblright{}.) Say it once more.
\end{tcolorbox}
\section[\texorpdfstring{\bn{দেনা} (denā)}{দেনা (denā)}]{\bn{দেনা} (denā) — a debt (owed); what one owes}

\label{lesson:BN-C416-dena}



\begin{quote}
Before the new one: say the Bengali for to lose (money or property), then the Bengali for an instalment.
\end{quote}

\subsection*{You'll want to know: \bn{দেনা}}
\textbf{\bn{দেনা}} — \emph{denā} — \textquotedblleft{}a debt (owed); what one owes\textquotedblright{}.

\begin{quote}
\bn{দোকানদারের} \bn{কাছে} \bn{আমাদের} \bn{কিছু} \bn{দেনা} \bn{আছে।}
\end{quote}

\begin{quote}
\emph{dokāndārer kāchhe āmāder kichhu denā āchhe.}
\end{quote}

\begin{quote}
We owe the shopkeeper some money.
\end{quote}

\bn{ধার} is a loan; \bn{দেনা} is what is still to be paid back.

\begin{cousinweb}
\textbf{\bn{দেনা}} (\emph{denā}) is built on \bn{দে}-, the stem of \bn{দেওয়া}, to give, the Bengali verb that goes back to Sanskrit \emph{dadāti} (Monier-Williams): what one has still to give.
\end{cousinweb}

\begin{grammarlens}[title={What you've built}]
One more word for getting around.
\end{grammarlens}

\subsection*{Your turn}
\begin{itemize}
\raggedright
  \item \emph{Say it:} \emph{denā}
  \item \emph{Say it:} \emph{denā}, once more
  \item \emph{Say it:} \emph{dokāndārer kāchhe āmāder kichhu denā āchhe.} — We owe the shopkeeper some money.
  \item \emph{Recall:} say the Bengali for to lose (money or property), then the Bengali for an instalment, then say \emph{denā} again
\end{itemize}

\begin{tcolorbox}[breakable,colback=teal!4,colframe=teal!35!black,title={Before you move on}]
What does \bn{দেনা} mean? (\textquotedblleft{}A debt (owed); what one owes\textquotedblright{}.) Say it once more.
\end{tcolorbox}
\section[\texorpdfstring{\bn{পাওনা} (pāonā)}{পাওনা (pāonā)}]{\bn{পাওনা} (pāonā) — money owed to one; one's dues}

\label{lesson:BN-C416-paona}



\begin{quote}
Before the new one: say the Bengali for an instalment, then the Bengali for a debt (owed).
\end{quote}

\subsection*{You'll want to know: \bn{পাওনা}}
\textbf{\bn{পাওনা}} — \emph{pāonā} — \textquotedblleft{}money owed to one; one's dues\textquotedblright{}.

\begin{quote}
\bn{কাজ} \bn{শেষ} \bn{করে} \bn{মিস্ত্রি} \bn{তার} \bn{পাওনা} \bn{চাইলেন।}
\end{quote}

\begin{quote}
\emph{kāj shesh kore mistri tār pāonā chāilen.}
\end{quote}

\begin{quote}
When the job was done, the mechanic asked for what he was owed.
\end{quote}

The pair \bn{দেনা} and \bn{পাওনা}, what you owe and what you are owed, often go together.

\begin{cousinweb}
\textbf{\bn{পাওনা}} (\emph{pāonā}) is built on \textbf{\bn{পাওয়া}}, \textquotedblleft{}to get\textquotedblright{}, from Sanskrit \emph{prāpayati} (CDIAL 8943): what one is due to get.
\end{cousinweb}

\begin{grammarlens}[title={What you've built}]
One more word for getting around.
\end{grammarlens}

\subsection*{Your turn}
\begin{itemize}
\raggedright
  \item \emph{Say it:} \emph{pāonā}
  \item \emph{Say it:} \emph{pāonā}, once more
  \item \emph{Say it:} \emph{kāj shesh kore mistri tār pāonā chāilen.} — When the job was done, the mechanic asked for what he was owed.
  \item \emph{Recall:} say the Bengali for an instalment, then the Bengali for a debt (owed), then say \emph{pāonā} again
\end{itemize}

\begin{tcolorbox}[breakable,colback=teal!4,colframe=teal!35!black,title={Before you move on}]
What does \bn{পাওনা} mean? (\textquotedblleft{}Money owed to one; one's dues\textquotedblright{}.) Say it once more.
\end{tcolorbox}
\section[\texorpdfstring{\bn{শোধ} \bn{করা} (shodh kôrā)}{শোধ করা (shodh kôrā)}]{\bn{শোধ} \bn{করা} (shodh kôrā) — to pay off; to repay (a debt)}

\label{lesson:BN-C416-shodh-kora}



\begin{quote}
Before the new one: say the Bengali for a debt (owed), then the Bengali for money owed to one.
\end{quote}

\subsection*{You'll want to know: \bn{শোধ} \bn{করা}}
\textbf{\bn{শোধ} \bn{করা}} — \emph{shodh kôrā} — \textquotedblleft{}to pay off; to repay (a debt)\textquotedblright{}.

\begin{quote}
\bn{দুই} \bn{বছরে} \bn{অরুণ} \bn{সব} \bn{ধার} \bn{শোধ} \bn{করল।}
\end{quote}

\begin{quote}
\emph{dui bôchhôre ôrun sôb dhār shodh kôrlô.}
\end{quote}

\begin{quote}
In two years Arun paid off all his debts.
\end{quote}

Past (he/she …): \bn{তিনি} \bn{ধার} \bn{শোধ} \bn{করলেন}, \emph{tini dhār shodh kôrlen}.

\begin{cousinweb}
\textbf{\bn{শোধ} \bn{করা}} (\emph{shodh kôrā}) joins Sanskrit \emph{śodha}, \textquotedblleft{}cleansing, payment\textquotedblright{} (Monier-Williams), to \bn{করা}, to do: to clear a debt. CDIAL lists the spoken Bengali \emph{sodhā}, \textquotedblleft{}to pay off\textquotedblright{}, from the same root (CDIAL 12626).
\end{cousinweb}

\begin{grammarlens}[title={What you've built}]
One more word for getting around.
\end{grammarlens}

\subsection*{Your turn}
\begin{itemize}
\raggedright
  \item \emph{Say it:} \emph{shodh kôrā}
  \item \emph{Say it:} \emph{shodh kôrā}, once more
  \item \emph{Say it:} \emph{dui bôchhôre ôrun sôb dhār shodh kôrlô.} — In two years Arun paid off all his debts.
  \item \emph{Recall:} say the Bengali for a debt (owed), then the Bengali for money owed to one, then say \emph{shodh kôrā} again
\end{itemize}

\begin{tcolorbox}[breakable,colback=teal!4,colframe=teal!35!black,title={Before you move on}]
What does \bn{শোধ} \bn{করা} mean? (\textquotedblleft{}To pay off; to repay (a debt)\textquotedblright{}.) Say it once more.
\end{tcolorbox}
\section[\texorpdfstring{\bn{মেটানো} (meṭāno)}{মেটানো (meṭāno)}]{\bn{মেটানো} (meṭāno) — to settle (a bill); to make good}

\label{lesson:BN-C416-metano}



\begin{quote}
Before the new one: say the Bengali for money owed to one, then the Bengali for to pay off.
\end{quote}

\subsection*{You'll want to know: \bn{মেটানো}}
\textbf{\bn{মেটানো}} — \emph{meṭāno} — \textquotedblleft{}to settle (a bill); to make good\textquotedblright{}.

\begin{quote}
\bn{হোটেল} \bn{ছাড়ার} \bn{আগে} \bn{আমরা} \bn{বিল} \bn{মেটালাম।}
\end{quote}

\begin{quote}
\emph{hoṭel chhāṛār āge āmrā bil meṭālām.}
\end{quote}

\begin{quote}
Before leaving the hotel we settled the bill.
\end{quote}

Past (he/she …): \bn{তিনি} \bn{বিল} \bn{মেটালেন}, \emph{tini bil meṭālen}.

Also of a quarrel or a need: \bn{ঝগড়া} \bn{মেটানো}, to settle a quarrel; \bn{খিদে} \bn{মেটানো}, to satisfy hunger.

\begin{cousinweb}
\textbf{\bn{মেটানো}} (\emph{meṭāno}) is the causing form of \bn{মেটা}, which CDIAL lists as Bengali \emph{meṭā}, \textquotedblleft{}to be wiped out\textquotedblright{}, under Sanskrit \emph{mṛṣṭá}, \textquotedblleft{}rubbed, washed\textquotedblright{} (CDIAL 10299): to settle a bill is to wipe it out.
\end{cousinweb}

\begin{grammarlens}[title={What you've built}]
One more word for getting around.
\end{grammarlens}

\subsection*{Your turn}
\begin{itemize}
\raggedright
  \item \emph{Say it:} \emph{meṭāno}
  \item \emph{Say it:} \emph{meṭāno}, once more
  \item \emph{Say it:} \emph{hoṭel chhāṛār āge āmrā bil meṭālām.} — Before leaving the hotel we settled the bill.
  \item \emph{Recall:} say the Bengali for money owed to one, then the Bengali for to pay off, then say \emph{meṭāno} again
\end{itemize}

\begin{tcolorbox}[breakable,colback=teal!4,colframe=teal!35!black,title={Before you move on}]
What does \bn{মেটানো} mean? (\textquotedblleft{}To settle (a bill); to make good\textquotedblright{}.) Say it once more.
\end{tcolorbox}
